\documentclass[tikz,border={24pt 24pt 24pt 60pt}]{standalone}
\usetikzlibrary{mindmap, trees, arrows.meta, shapes.geometric} % <-- Added arrows and shapes
\usepackage{fontawesome5}

% --- NAGWA BRAND COLORS ---
% Text/Line Colors
\definecolor{Grey}{RGB}{76,76,78}
\definecolor{Brown}{RGB}{138,93,59}
\definecolor{Blue}{RGB}{85,164,224}
\definecolor{BBlue}{RGB}{0,115,181}
\definecolor{Red}{RGB}{225,75,59}
\definecolor{Orange}{RGB}{243,162,3}
\definecolor{Yellow}{RGB}{252,238,33}
\definecolor{Lemon}{RGB}{203,219,67}
\definecolor{Green}{RGB}{16,158,36}
\definecolor{GGreen}{RGB}{14,127,71}
\definecolor{GGGreen}{RGB}{0,156,156}
\definecolor{Purple}{RGB}{124,36,231}
\definecolor{PPurple}{RGB}{163,10,120}
\definecolor{Fushia}{RGB}{225,50,162}

% Fill Colors
\definecolor{RootGrey}{HTML}{E0E0E0} % <-- Replace E0E0E0 with your exact grey pastel hex code
\definecolor{Pink1}{HTML}{F4D4E1}
\definecolor{Pink2}{HTML}{F7E2EB}
\definecolor{Peach1}{HTML}{F6C9A0}
\definecolor{Peach2}{HTML}{FCDDC5}
\definecolor{Olive1}{HTML}{CED5B2}
\definecolor{Olive2}{HTML}{E6EAD6}
\definecolor{Sage1}{HTML}{AFC3BC}
\definecolor{Sage2}{HTML}{D3E0DB}
\definecolor{Blue1}{HTML}{ACC5DB}
\definecolor{Blue2}{HTML}{D3E0EA}
\definecolor{Purple1}{HTML}{C6A5C0}
\definecolor{Purple2}{HTML}{E8D1E4}
\definecolor{Rose1}{HTML}{FF9FAF}
\definecolor{Rose2}{HTML}{FFD4DD}

% --- BILINGUAL SETUP ---
% Fonts are chosen automatically by language:
%   Arabic → Noto (Regular / Bold)
%   English → STIX2Text family
\usepackage{polyglossia}
<% if lang == 'ar' -%>
\setmainlanguage{arabic}
\setotherlanguage{english}
\newfontfamily\arabicfont[
  Path=<< font_dir >>,
  Script=Arabic,
  BoldFont=<< arabic_font_bold_file >>,
]{<< arabic_font_file >>}
\newfontfamily\englishfont[
  Path=<< font_dir >>,
  BoldFont=<< english_font_bold_file >>,
  ItalicFont=<< english_font_italic_file >>,
  BoldItalicFont=<< english_font_bolditalic_file >>,
]{<< english_font_file >>}
<% else -%>
\setmainlanguage{english}
\setotherlanguage{arabic}
\setmainfont[
  Path=<< font_dir >>,
  BoldFont=<< english_font_bold_file >>,
  ItalicFont=<< english_font_italic_file >>,
  BoldItalicFont=<< english_font_bolditalic_file >>,
]{<< english_font_file >>}
% Arabic font is always loaded so RTL snippets can appear in English maps.
\newfontfamily\arabicfont[
  Path=<< font_dir >>,
  Script=Arabic,
  BoldFont=<< arabic_font_bold_file >>,
]{<< arabic_font_file >>}
<% endif -%>

\begin{document}

<% set text_dir = '\setRTL' if lang == 'ar' else '\setLTR' -%>
<% set n = root.children|length -%>

<% set ns = namespace(has_level_3=false) -%>
<% for child in root.children -%>
  <% if child.children|length > 0 -%>
    <% set ns.has_level_3 = true -%>
  <% endif -%>
<% endfor -%>

\begin{tikzpicture}[
<% if layout_mode == 'radial' -%>
    grow cyclic,
    mindmap,
    <% if shape == 'rectangle' -%>
    every node/.style={rectangle, rounded corners=3pt, inner xsep=10pt, inner ysep=10pt, align=center, execute at begin node={<< text_dir >>\hyphenpenalty=10000}},
    every child/.append style={
        edge from parent path={(\tikzparentnode) -- (\tikzchildnode)},
        edge from parent/.style={draw, black!50, thin}
    },
    color logic child/.style={},
    color logic node/.style={fill=#1},
    <% else -%>
    every node/.style={concept, inner sep=3pt, align=center, execute at begin node={<< text_dir >>\hyphenpenalty=10000}},
    <% if child_path == 'grey' -%>
    concept color=black!10,
    color logic child/.style={},
    color logic node/.style={fill=#1},
    <% else -%>
    concept color=<< root.fill >>,
    color logic child/.style={concept color=#1},
    color logic node/.style={},
    <% endif -%>
    <% endif -%>
    level 1/.append style={
        level distance=7cm,
        sibling angle=<< 360 // n if n > 0 else 360 >>,
        <% if not ns.has_level_3 and n % 2 != 0 -%>
        <% if lang == 'ar' -%>counterclockwise from=90,<% else -%>clockwise from=90,<% endif -%>
        <% endif -%>
        text width=3cm,
        font=\fontsize{13.5pt}{16.2pt}\selectfont\bfseries,
        inner sep=10pt,
        minimum size=2.25cm
    },
    level 2/.append style={
        level distance=6cm,
        sibling angle=45,
        text width=3.5cm,
        font=\fontsize{12pt}{14.4pt}\selectfont\bfseries,
        inner sep=10pt,
        minimum size=1.75cm
    },
    level 3/.append style={
        level distance=6cm,
        sibling angle=90,
        text width=2.5cm,
        font=\fontsize{12pt}{14.4pt}\selectfont\bfseries,
        inner sep=10pt,
        minimum size=1.15cm
    }
<% else -%>
    % --- VERTICAL & HORIZONTAL TREES ---
    <% if layout_mode == 'vertical' -%>
    grow=down,
    level 1/.style={sibling distance=6.5cm, level distance=3cm},
    level 2/.style={sibling distance=3.5cm, level distance=3cm},
    level 3/.style={sibling distance=2cm, level distance=3cm},
    <% elif layout_mode == 'horizontal' -%>
    <% if lang == 'ar' %>grow=left,<% else %>grow=right,<% endif %>
    level 1/.style={sibling distance=4cm, level distance=5.5cm},
    level 2/.style={sibling distance=2.5cm, level distance=5cm},
    level 3/.style={sibling distance=1.5cm, level distance=4.5cm},
    <% endif -%>
    % Node Styles for Trees
    every node/.style={
        <% if shape == 'rectangle' %>rectangle, rounded corners=3pt,<% else %>circle,<% endif %>
        text width=3.5cm,
        font=\fontsize{12pt}{14.4pt}\selectfont\bfseries,
        inner xsep=10pt, inner ysep=10pt, align=center,
        execute at begin node={<< text_dir >>\hyphenpenalty=10000}
    },
    color logic node/.style={fill=#1, draw=black!20, thick},
    color logic child/.style={},
    % Edge Styles for Trees
    <% if edge_style == 'orthogonal' -%>
        <% if layout_mode == 'vertical' -%>
        edge from parent path={(\tikzparentnode.south) -- +(0,-0.5) -| (\tikzchildnode.north)},
        <% else -%>
        <% if lang == 'ar' %>
        edge from parent path={(\tikzparentnode.west) -- +(-0.5,0) |- (\tikzchildnode.east)},
        <% else %>
        edge from parent path={(\tikzparentnode.east) -- +(0.5,0) |- (\tikzchildnode.west)},
        <% endif %>
        <% endif -%>
        edge from parent/.style={draw, black!60, thick},
    <% elif edge_style == 'arrows' -%>
        edge from parent path={(\tikzparentnode) -- (\tikzchildnode)},
        edge from parent/.style={draw, black!60, thick, -{Stealth[scale=1.2]}},
    <% else -%>
        % Standard smooth curve
        <% if layout_mode == 'vertical' -%>
        edge from parent path={(\tikzparentnode.south) .. controls +(0,-1.5) and +(0,1.5) .. (\tikzchildnode.north)},
        <% else -%>
        <% if lang == 'ar' %>
        edge from parent path={(\tikzparentnode.west) .. controls +(-1.5,0) and +(1.5,0) .. (\tikzchildnode.east)},
        <% else %>
        edge from parent path={(\tikzparentnode.east) .. controls +(1.5,0) and +(-1.5,0) .. (\tikzchildnode.west)},
        <% endif %>
        <% endif -%>
        edge from parent/.style={draw, black!60, thick},
    <% endif -%>
<% endif -%>
]

  \node[color logic node=<< root.fill >> ... % (Keep your existing node loop logic from here down)

  \node[color logic node=<< root.fill >>, text=<< root.text_color >>, font=\fontsize{13.5pt}{16.2pt}\selectfont\bfseries, inner sep=10pt, minimum size=2.25cm<% if root.custom %>, << root.custom >><% endif %>] {
    <% if root.icon %>\LR{\huge\faIcon{<< root.icon >>}}\\[0.5ex]<% endif %><< root.text >>
  }
    <% for child in root.children %>
      child[color logic child=<< child.fill >><% if child.custom %>, << child.custom >><% endif %>] { node[color logic node=<< child.fill >>, text=<< child.text_color >><% if child.custom %>, << child.custom >><% endif %>] {
        <% if child.icon %>\LR{\huge\faIcon{<< child.icon >>}}\\[0.5ex]<% endif %><< child.text >>
      }
        <% if child.children %>
          <% for grandchild in child.children %>
            child[color logic child=<< grandchild.fill >><% if grandchild.custom %>, << grandchild.custom >><% endif %>] { node[color logic node=<< grandchild.fill >>, text=<< grandchild.text_color >><% if grandchild.custom %>, << grandchild.custom >><% endif %>] {
              <% if grandchild.icon %>\LR{\LARGE\faIcon{<< grandchild.icon >>}}\\[0.5ex]<% endif %><< grandchild.text >>
            } }
          <% endfor %>
        <% endif %>
      }
    <% endfor %>
  ;
\end{tikzpicture}
\end{document}
